%==============================================================================
%  Standalone note: the velocity-magnitude derivative floor.
%
%  Documents the single numerical regularization parameter
%      velocity_magnitude_derivative_floor      (Julia: PhysicalProperties field;
%                                                 default 1e-12, base_config.json)
%  that keeps the Exact-Newton Jacobian finite at a stagnation point.
%
%  Notation is kept compatible with the manuscript
%  "A stabilized finite element method for incompressible, inertial flows
%   in inhomogeneous porous media" (Casas, Gonzalez-Usua, Codina, de-Pouplana):
%      sigma(alpha,u) = a(alpha) + b(alpha)|u|         (eq:DBFResistanceTerm)
%      tau_1, tau_2, tau_{1,NS}                          (eq:Tau1, eq:Tau2, eq:TauNavierStokes)
%==============================================================================
\documentclass[11pt]{article}

\usepackage[utf8]{inputenc}
\usepackage[margin=1in]{geometry}
\usepackage{amsmath,amssymb,amsthm,mathtools}
\usepackage{bm}
\usepackage[colorlinks=true,allcolors=blue]{hyperref}

\newcommand{\eps}{\varepsilon_{\mathrm{d}}}          % the derivative floor
\newcommand{\uu}{\bm{u}}
\newcommand{\du}{\mathrm{d}\bm{u}}
\newcommand{\norm}[1]{\lvert #1 \rvert}

\title{The velocity-magnitude derivative floor\\
       \large Regularizing $\partial\norm{\uu}/\partial\uu$ in the Exact-Newton Jacobian}
\author{Porous Navier--Stokes solver --- implementation note}
\date{}

\begin{document}
\maketitle

\noindent
\textbf{Code anchor.} Julia parameter \texttt{PhysicalProperties.velocity\_magnitude\_derivative\_floor}
(default $\eps = 10^{-12}$ in \texttt{config/base\_config.json}); carried on the velocity-regularization
policy \texttt{SmoothVelocityFloor} and consumed at three leaf sites:
\texttt{dsigma\_du} (\texttt{src/models/reaction.jl}),
\texttt{compute\_dtau\_1\_du} and \texttt{compute\_dtau\_2\_du} (\texttt{src/stabilization/tau.jl}).

\section{Where it arises}

Both the reaction coefficient and the stabilization parameters depend on the velocity \emph{only}
through the Euclidean speed
\begin{equation}
  \norm{\uu} \;=\; \sqrt{\uu\cdot\uu}\,,
\end{equation}
namely
\begin{equation}
  \sigma(\alpha,\uu) = a(\alpha) + b(\alpha)\,\norm{\uu},
  \qquad
  \tau_{1,\mathrm{NS}}^{-1} = \frac{c_1\nu}{h^2} + \frac{c_2\,\norm{\uu}}{h},
\end{equation}
with $\tau_1=\big(\alpha\,\tau_{1,\mathrm{NS}}^{-1}+\sigma\big)^{-1}$ and
$\tau_2 = h^2/\!\big(c_1\alpha\,\tau_{1,\mathrm{NS}}\big)$.

The Exact-Newton tangent requires the directional derivatives $\partial\sigma/\partial\uu$,
$\partial\tau_1/\partial\uu$ and $\partial\tau_2/\partial\uu$, every one of which factors through the
gradient of the speed,
\begin{equation}\label{eq:normderiv}
  \frac{\partial \norm{\uu}}{\partial \uu}\cdot\du
  \;=\; \frac{\uu\cdot\du}{\norm{\uu}}\,.
\end{equation}

\section{The singularity and its regularization}

The map $\uu\mapsto\norm{\uu}$ is only $C^0$ at the origin: \eqref{eq:normderiv} is the indeterminate
$0/0$ at $\uu=\bm 0$. A stagnation point with $\uu=\bm 0$ is not exotic --- it occurs exactly on
perfect Dirichlet walls (and at interior stagnation points of a manufactured field), where the discrete
velocity is pinned to zero. Evaluated verbatim, \eqref{eq:normderiv} returns \texttt{NaN} and poisons the
whole Jacobian.

We regularize the denominator with a single additive floor $\eps>0$:
\begin{equation}\label{eq:reg}
  \frac{\partial \norm{\uu}}{\partial \uu}\cdot\du
  \;\approx\; \frac{\uu\cdot\du}{\norm{\uu}+\eps}\,.
\end{equation}
This is the only role of \texttt{velocity\_magnitude\_derivative\_floor}. Concretely it enters
\begin{align}
  \partial_{\uu}\sigma\cdot\du
    &= b(\alpha)\,\frac{\uu\cdot\du}{\norm{\uu}+\eps}
    && \text{(\texttt{dsigma\_du}, Forchheimer)},\\
  \partial_{\uu}\tau_{1,\mathrm{NS}}^{-1}\cdot\du
    &= \frac{c_2}{h}\,\frac{\uu\cdot\du}{\norm{\uu}+\eps}
    && \text{(inside \texttt{compute\_dtau\_1\_du}, \texttt{compute\_dtau\_2\_du})},
\end{align}
with $\partial_{\uu}\tau_1=-\tau_1^2\big(\alpha\,\partial_{\uu}\tau_{1,\mathrm{NS}}^{-1}
+\partial_{\uu}\sigma\big)$ and the analogous chain for $\tau_2$.

\section{Why it cannot move the converged solution}

Three facts pin down the status of $\eps$ as a \emph{Jacobian-only} regularization:
\begin{enumerate}
  \item \textbf{It is absent from the residual.} The nonlinear residual $\bm R(\uu,p)$ that the solver
        drives to zero contains $\sigma\,\uu$, $\tau_1$, $\tau_2$ --- but never \eqref{eq:reg}. The fixed
        point $\bm R=\bm 0$ is therefore independent of $\eps$.
  \item \textbf{It is dropped in Picard.} In \texttt{PicardMode} the coefficient derivatives
        $\partial_{\uu}\sigma$, $\partial_{\uu}\tau$ are frozen to structural zeros, so $\eps$ does not
        appear at all; the Picard tangent is exact in $\eps$ trivially.
  \item \textbf{It is an inexact-Newton perturbation.} In \texttt{ExactNewtonMode} it perturbs the tangent
        by $O\!\big(\eps/\norm{\uu}\big)$, which vanishes in the bulk where $\norm{\uu}\gg\eps$ and is
        bounded near stagnation. A perturbed tangent changes only the Newton \emph{convergence rate}, never
        the root it converges to.
\end{enumerate}
Hence $\eps$ is a pure conditioning knob: it guarantees a finite, well-defined Jacobian everywhere while
leaving the discrete solution --- and thus every reported error and convergence rate --- untouched.

\section{Interaction with the speed floors (no double counting)}

The speed handed to \eqref{eq:reg} is already the \emph{regularized} speed produced by the velocity policy:
\texttt{reaction\_speed} $=\sqrt{\uu\cdot\uu+u_{\mathrm{base}}^2}$ for $\sigma$, and
\texttt{effective\_speed} (which additionally carries a mesh-diffusive floor) for $\tau$. Thus whenever the
base floor $u_{\mathrm{base}}=\texttt{u\_base\_floor\_ref}>0$ one already has $\norm{\uu}\ge u_{\mathrm{base}}\gg\eps$
and \eqref{eq:reg} is a no-op. The floor $\eps$ is load-bearing precisely in the regime the convergence
studies use, where the harness sets $u_{\mathrm{base}}=0$ to keep the drag law mesh-independent: there
$\norm{\uu}$ can reach $0$ and $\eps$ is the only guard against the $0/0$. Keeping $\eps$ a \emph{separate}
parameter from \texttt{u\_base\_floor\_ref} and \texttt{epsilon\_floor} (which guards an unrelated
$h\to0$ denominator) lets the three be reasoned about and tuned independently.

\section{Recommended value}

$\eps=10^{-12}$ (double precision near the rounding floor) is small enough to be numerically invisible
wherever $\norm{\uu}$ is not itself at the level of machine epsilon, yet large enough to keep
$\uu/(\norm{\uu}+\eps)$ finite. Because the parameter only conditions the Jacobian, any value in a wide
band $\eps\in[10^{-14},10^{-8}]$ gives indistinguishable solves; it is exposed as a config input purely so
the value is documented, validated at load, and tunable rather than buried as a literal.

\end{document}
